\section*{Заключение}
\addcontentsline{toc}{section}{Заключение}

В работе проведено сравнительное исследование 10 нейросетевых архитектур на задаче NER в научных текстах (датасет SciERC).

\textbf{Основные результаты.} Лучший результат в базовом эксперименте показала модель \texttt{scibert\_linear\_noweight} --- macro F1 = 0.646. SciBERT стабильно превосходит BERT и RoBERTa за счёт научного предобучения. Простая линейная голова без взвешивания классов оказалась конкурентоспособнее сложных вариантов (MLP, CRF, Concat-4).

\textbf{Сильные стороны.}
\begin{itemize}
    \item Контролируемое сравнение: все модели обучены с одинаковыми гиперпараметрами, единственная переменная --- архитектура.
    \item Полный пайплайн: от данных до confusion matrix и кривых обучения.
    \item Воспроизводимость: фиксированный seed, реестр прогонов.
\end{itemize}

\textbf{Слабые стороны.}
\begin{itemize}
    \item Одна случайная инициализация на модель --- нет доверительных интервалов.
    \item CRF-модели дали идентичные результаты в двух вариантах (вероятно, технический артефакт Kaggle-сессии).
    \item Редкие типы (Metric, Task) остаются сложными для всех моделей.
\end{itemize}

\textbf{Направления развития.}
\begin{itemize}
    \item Span-based модели (DyGIE++, SpanBERT) вместо token-level BIO --- качественно другой подход к границам сущностей.
    \item Joint NER+RE обучение, поскольку SciERC содержит разметку отношений.
    \item Более крупные энкодеры (SciBERT-large, SPECTER).
\end{itemize}
